% GENERATED FILE. Edit canonical lessons, then run npm run generate:books.
% canonical-source-hash: fnv1a64:456c60139b1a2638
% canonical-lessons: TA-C418-vituti, TA-C418-varaverpalar, TA-C418-mintukki, TA-C418-patukkai, TA-C418-viri

\chapter{At the Lodge}
\label{ch:ta-at-the-lodge}
\hlchaptermodality{418}

\begin{chapteropening}
\textbf{By the end of this chapter:} \emph{I can talk about a lodge, a receptionist, a lift and a bed, and spreading a sheet.}

Complete the last lesson of chapter 418: I can talk about a lodge, a receptionist, a lift and a bed, and spreading a sheet.
\end{chapteropening}

\section[\texorpdfstring{\ta{விடுதி} (viṭuti)}{விடுதி (viṭuti)}]{\ta{விடுதி} (viṭuti) — a lodge; a hostel}

\label{lesson:TA-C418-vituti}



\begin{quote}
Before the new one: say the Tamil for to flounder, then the Tamil for a shore.
\end{quote}

\subsection*{You'll want to know: \ta{விடுதி}}
\textbf{\ta{விடுதி}} — \emph{viṭuti} — \textquotedblleft{}a lodge; a hostel\textquotedblright{}.

\begin{quote}
\ta{நாங்கள்} \ta{ஒரு} \ta{சிறிய} \ta{விடுதியில்} \ta{தங்கினோம்}.
\end{quote}

\begin{quote}
\emph{nāṅkaḷ oru ciṟiya viṭutiyil taṅkiṉōm.}
\end{quote}

\begin{quote}
We stayed in a small lodge.
\end{quote}

\ta{மாணவர்} \ta{விடுதி} is a students' hostel.

\begin{cousinweb}
\textbf{\ta{விடுதி}} (\emph{viṭuti}) is native Dravidian: DEDR 5393 lists it under \ta{விடு}, \textquotedblleft{}to leave, to let go\textquotedblright{}, as \textquotedblleft{}a lodging\textquotedblright{}.
\end{cousinweb}

\begin{grammarlens}[title={What you've built}]
One more word for getting around.
\end{grammarlens}

\subsection*{Your turn}
\begin{itemize}
\raggedright
  \item \emph{Say it:} \emph{viṭuti}
  \item \emph{Say it:} \emph{viṭuti}, once more
  \item \emph{Say it:} \emph{nāṅkaḷ oru ciṟiya viṭutiyil taṅkiṉōm.} — We stayed in a small lodge.
  \item \emph{Recall:} say the Tamil for to flounder, then the Tamil for a shore, then say \emph{viṭuti} again
\end{itemize}

\begin{tcolorbox}[breakable,colback=teal!4,colframe=teal!35!black,title={Before you move on}]
What does \ta{விடுதி} mean? (\textquotedblleft{}A lodge; a hostel\textquotedblright{}.) Say it once more.
\end{tcolorbox}
\section[\texorpdfstring{\ta{வரவேற்பாளர்} (varavēṟpāḷar)}{வரவேற்பாளர் (varavēṟpāḷar)}]{\ta{வரவேற்பாளர்} (varavēṟpāḷar) — a receptionist}

\label{lesson:TA-C418-varaverpalar}



\begin{quote}
Before the new one: say the Tamil for a shore, then the Tamil for a lodge.
\end{quote}

\subsection*{You'll want to know: \ta{வரவேற்பாளர்}}
\textbf{\ta{வரவேற்பாளர்}} — \emph{varavēṟpāḷar} — \textquotedblleft{}a receptionist\textquotedblright{}.

\begin{quote}
\ta{வரவேற்பாளர்} \ta{எங்களுக்கு} \ta{அறைச்} \ta{சாவியைக்} \ta{கொடுத்தார்}.
\end{quote}

\begin{quote}
\emph{varavēṟpāḷar eṅkaḷukku aṟaic cāviyaik koṭuttār.}
\end{quote}

\begin{quote}
The receptionist gave us the room key.
\end{quote}

The ending -\ta{ஆளர்} makes a person: \ta{வரவேற்பு}, a welcome, gives \ta{வரவேற்பாளர்}.

\begin{cousinweb}
\textbf{\ta{வரவேற்பாளர்}} (\emph{varavēṟpāḷar}) is built from \ta{வரவேற்பு}, \textquotedblleft{}a welcome\textquotedblright{} (\ta{வரவு}, \textquotedblleft{}coming\textquotedblright{}, DEDR 5270, + \ta{ஏற்பு}, \textquotedblleft{}receiving\textquotedblright{}, DEDR 905), with the ending \emph{-āḷar}, \textquotedblleft{}one who\textquotedblright{}.
\end{cousinweb}

\begin{grammarlens}[title={What you've built}]
One more word for getting around.
\end{grammarlens}

\subsection*{Your turn}
\begin{itemize}
\raggedright
  \item \emph{Say it:} \emph{varavēṟpāḷar}
  \item \emph{Say it:} \emph{varavēṟpāḷar}, once more
  \item \emph{Say it:} \emph{varavēṟpāḷar eṅkaḷukku aṟaic cāviyaik koṭuttār.} — The receptionist gave us the room key.
  \item \emph{Recall:} say the Tamil for a shore, then the Tamil for a lodge, then say \emph{varavēṟpāḷar} again
\end{itemize}

\begin{tcolorbox}[breakable,colback=teal!4,colframe=teal!35!black,title={Before you move on}]
What does \ta{வரவேற்பாளர்} mean? (\textquotedblleft{}A receptionist\textquotedblright{}.) Say it once more.
\end{tcolorbox}
\section[\texorpdfstring{\ta{மின்தூக்கி} (miṉtūkki)}{மின்தூக்கி (miṉtūkki)}]{\ta{மின்தூக்கி} (miṉtūkki) — an elevator; a lift}

\label{lesson:TA-C418-mintukki}



\begin{quote}
Before the new one: say the Tamil for a lodge, then the Tamil for a receptionist.
\end{quote}

\subsection*{You'll want to know: \ta{மின்தூக்கி}}
\textbf{\ta{மின்தூக்கி}} — \emph{miṉtūkki} — \textquotedblleft{}an elevator; a lift\textquotedblright{}.

\begin{quote}
\ta{மின்தூக்கி} \ta{வேலை} \ta{செய்யாததால்} \ta{படிக்கட்டில்} \ta{ஏறினோம்}.
\end{quote}

\begin{quote}
\emph{miṉtūkki vēlai ceyyātatāl paṭikkaṭṭil ēṟiṉōm.}
\end{quote}

\begin{quote}
Because the lift was not working, we climbed the stairs.
\end{quote}

\ta{மின்}, electric, + \ta{தூக்கி}, a lifter.

\begin{cousinweb}
\textbf{\ta{மின்தூக்கி}} (\emph{miṉtūkki}) is a modern coinage built from \ta{மின்}, \textquotedblleft{}lightning, electric\textquotedblright{} (DEDR 4876), and \ta{தூக்கி}, \textquotedblleft{}a lifter\textquotedblright{}, from \ta{தூக்கு}, \textquotedblleft{}to lift\textquotedblright{} (DEDR 3376).
\end{cousinweb}

\begin{grammarlens}[title={What you've built}]
One more word for getting around.
\end{grammarlens}

\subsection*{Your turn}
\begin{itemize}
\raggedright
  \item \emph{Say it:} \emph{miṉtūkki}
  \item \emph{Say it:} \emph{miṉtūkki}, once more
  \item \emph{Say it:} \emph{miṉtūkki vēlai ceyyātatāl paṭikkaṭṭil ēṟiṉōm.} — Because the lift was not working, we climbed the stairs.
  \item \emph{Recall:} say the Tamil for a lodge, then the Tamil for a receptionist, then say \emph{miṉtūkki} again
\end{itemize}

\begin{tcolorbox}[breakable,colback=teal!4,colframe=teal!35!black,title={Before you move on}]
What does \ta{மின்தூக்கி} mean? (\textquotedblleft{}An elevator; a lift\textquotedblright{}.) Say it once more.
\end{tcolorbox}
\section[\texorpdfstring{\ta{படுக்கை} (paṭukkai)}{படுக்கை (paṭukkai)}]{\ta{படுக்கை} (paṭukkai) — bedding; a bed}

\label{lesson:TA-C418-patukkai}



\begin{quote}
Before the new one: say the Tamil for a receptionist, then the Tamil for an elevator.
\end{quote}

\subsection*{You'll want to know: \ta{படுக்கை}}
\textbf{\ta{படுக்கை}} — \emph{paṭukkai} — \textquotedblleft{}bedding; a bed\textquotedblright{}.

\begin{quote}
\ta{அறையில்} \ta{இரண்டு} \ta{படுக்கைகள்} \ta{இருந்தன}.
\end{quote}

\begin{quote}
\emph{aṟaiyil iraṇṭu paṭukkaikaḷ iruntaṉa.}
\end{quote}

\begin{quote}
There were two beds in the room.
\end{quote}

\ta{படுக்கையறை}, a bedroom, is built on it.

\begin{cousinweb}
\textbf{\ta{படுக்கை}} (\emph{paṭukkai}) is native Dravidian (DEDR 3852), from \ta{படு}, \textquotedblleft{}to lie down\textquotedblright{}.
\end{cousinweb}

\begin{grammarlens}[title={What you've built}]
One more word for getting around.
\end{grammarlens}

\subsection*{Your turn}
\begin{itemize}
\raggedright
  \item \emph{Say it:} \emph{paṭukkai}
  \item \emph{Say it:} \emph{paṭukkai}, once more
  \item \emph{Say it:} \emph{aṟaiyil iraṇṭu paṭukkaikaḷ iruntaṉa.} — There were two beds in the room.
  \item \emph{Recall:} say the Tamil for a receptionist, then the Tamil for an elevator, then say \emph{paṭukkai} again
\end{itemize}

\begin{tcolorbox}[breakable,colback=teal!4,colframe=teal!35!black,title={Before you move on}]
What does \ta{படுக்கை} mean? (\textquotedblleft{}Bedding; a bed\textquotedblright{}.) Say it once more.
\end{tcolorbox}
\section[\texorpdfstring{\ta{விரி} (viri)}{விரி (viri)}]{\ta{விரி} (viri) — to spread out; to unfold}

\label{lesson:TA-C418-viri}



\begin{quote}
Before the new one: say the Tamil for an elevator, then the Tamil for bedding.
\end{quote}

\subsection*{You'll want to know: \ta{விரி}}
\textbf{\ta{விரி}} — \emph{viri} — \textquotedblleft{}to spread out; to unfold\textquotedblright{}.

\begin{quote}
\ta{அம்மா} \ta{படுக்கையில்} \ta{புதிய} \ta{போர்வையை} \ta{விரித்தாள்}.
\end{quote}

\begin{quote}
\emph{ammā paṭukkaiyil putiya pōrvaiyai virittāḷ.}
\end{quote}

\begin{quote}
Mother spread a new blanket on the bed.
\end{quote}

Past (I …): \ta{விரித்தேன்}, \emph{virittēṉ}.

\ta{விரித்தேன்} is I spread something; \ta{விரிந்தது}, it opened out.

\begin{cousinweb}
\textbf{\ta{விரி}} (\emph{viri}) is native Dravidian (DEDR 5411), \textquotedblleft{}to expand, to spread out, to unfold\textquotedblright{}.
\end{cousinweb}

\begin{grammarlens}[title={What you've built}]
One more word for getting around.
\end{grammarlens}

\subsection*{Your turn}
\begin{itemize}
\raggedright
  \item \emph{Say it:} \emph{viri}
  \item \emph{Say it:} \emph{viri}, once more
  \item \emph{Say it:} \emph{ammā paṭukkaiyil putiya pōrvaiyai virittāḷ.} — Mother spread a new blanket on the bed.
  \item \emph{Recall:} say the Tamil for an elevator, then the Tamil for bedding, then say \emph{viri} again
\end{itemize}

\begin{tcolorbox}[breakable,colback=teal!4,colframe=teal!35!black,title={Before you move on}]
What does \ta{விரி} mean? (\textquotedblleft{}To spread out; to unfold\textquotedblright{}.) Say it once more.
\end{tcolorbox}
